\begin{thebibliography}{18}
\bibitem[1]{ref1} Alexander E. Andreev. On a method for obtaining more than quadratic effective lower bounds for the complexity of $\pi$-schemes. Moscow University Mathematics Bulletin, 42(1):24–29, 1987.
\bibitem[2]{ref2} Susanna F. de Rezende, Or Meir, Jakob Nordström, Toniann Pitassi, and Robert Robere. KRW composition theorems via lifting. In Sandy Irani, editor, 61st IEEE Annual Symposium on Foundations of Computer Science, FOCS 2020, Durham, NC, USA, November 16-19, 2020, pages 43–49. IEEE, 2020. doi:10.1109/FOCS46700.2020.00013.
\bibitem[3]{ref3} Irit Dinur and Or Meir. Toward the KRW composition conjecture: Cubic formula lower bounds via communication complexity. In Ran Raz, editor, 31st Conference on Computational Complexity, CCC 2016, May 29 to June 1, 2016, Tokyo, Japan, volume 50 of LIPIcs, pages 3:1–3:51. Schloss Dagstuhl - Leibniz-Zentrum für Informatik, 2016. doi:10.4230/LIPIcs.CCC.2016.3.
\bibitem[4]{ref4} Jeff Edmonds, Russell Impagliazzo, Steven Rudich, and Jirí Sgall. Communication complexity towards lower bounds on circuit depth. Comput. Complex., 10(3):210–246, 2001. doi:10.1007/s00037-001-8195-x.
\bibitem[5]{ref5} Dmitry Gavinsky, Or Meir, Omri Weinstein, and Avi Wigderson. Toward better formula lower bounds: The composition of a function and a universal relation. SIAM J. Comput., 46(1):114–131, 2017. doi:10.1137/15M1018319.
\bibitem[6]{ref6} Johan Håstad. The shrinkage exponent of de morgan formulas is 2. SIAM J. Comput., 27(1):48–64, 1998. doi:10.1137/S0097539794261556.
\bibitem[7]{ref7} Johan Håstad and Avi Wigderson. Composition of the universal relation. In Advances In Computational Complexity Theory, Proceedings of a DIMACS Workshop, New Jersey, USA, December 3-7, 1990, pages 119–134, 1990. doi:10.1090/dimacs/013/07.
\bibitem[8]{ref8} Russell Impagliazzo, Raghu Meka, and David Zuckerman. Pseudorandomness from shrinkage. In 2012 IEEE 53rd Annual Symposium on Foundations of Computer Science, pages 111–119. IEEE, 2012.
\bibitem[9]{ref9} Russell Impagliazzo and Noam Nisan. The effect of random restrictions on formula size. Random Struct. Algorithms, 4:121–134, 1993.
\bibitem[10]{ref10} Mauricio Karchmer, Ran Raz, and Avi Wigderson. Super-logarithmic depth lower bounds via the direct sum in communication complexity. Computational Complexity, 5(3/4):191–204, 1995. doi:10.1007/BF01206317.
\bibitem[11]{ref11} Mauricio Karchmer and Avi Wigderson. Monotone circuits for connectivity require super- logarithmic depth. In Proceedings of the 20th Annual ACM Symposium on Theory of Computing, May 2-4, 1988, Chicago, Illinois, USA, pages 539–550, 1988. doi:10.1145/62212.62265.
\bibitem[12]{ref12} Valeriy Mihailovich Khrapchenko. Complexity of the realization of a linear function in the class of II-circuits. Mathematical Notes of the Academy of Sciences of the USSR, 9(1):21–23, 1971.
\bibitem[13]{ref13} Sajin Koroth and Or Meir. Improved composition theorems for functions and relations. In Eric Blais, Klaus Jansen, José D. P. Rolim, and David Steurer, editors, Approximation, Randomiza- tion, and Combinatorial Optimization. Algorithms and Techniques, APPROX/RANDOM 2018, August 20-22, 2018 - Princeton, NJ, USA, volume 116 of LIPIcs, pages 48:1–48:18. Schloss Dag- stuhl - Leibniz-Zentrum für Informatik, 2018. doi:10.4230/LIPIcs.APPROX-RANDOM.2018.48.
\bibitem[14]{ref14} Ivan Mihajlin and Alexander Smal. Toward better depth lower bounds: The XOR-KRW conjecture. In Valentine Kabanets, editor, 36th Computational Complexity Conference, CCC 2021, July 20-23, 2021, Toronto, Ontario, Canada (Virtual Conference), volume 200 of LIPIcs, pages 38:1–38:24. Schloss Dagstuhl - Leibniz-Zentrum für Informatik, 2021. doi:10.4230/LIPIcs.CCC.2021.38.
\bibitem[15]{ref15} Michael S. Paterson and Uri Zwick. Shrinkage of de morgan formulae under restriction. In Proceedings of the 32nd Annual Symposium on Foundations of Computer Science, SFCS ’91, pages 324–333, USA, 1991. IEEE Computer Society. doi:10.1109/SFCS.1991.185385.
\bibitem[16]{ref16} J. Riordan and C. Shannon. The number of two-terminal series-parallel networks. J Math Phys, 21, January 1942. doi:10.1002/sapm194221183.
\bibitem[17]{ref17} Bella Abramovna Subbotovskaya. Realization of linear functions by formulas using $\wedge$, $\vee$, $\neg$. In Doklady Akademii Nauk, volume 136-3, pages 553–555. Russian Academy of Sciences, 1961.
\bibitem[18]{ref18} Avishay Tal. Shrinkage of de morgan formulae by spectral techniques. Proceedings - Annual IEEE Symposium on Foundations of Computer Science, FOCS, pages 551–560, December 2014. doi:10.1109/FOCS.2014.65.
\end{thebibliography}
